\documentclass[11pt,usenames,dvipsnames]{beamer}

\usetheme[numbering=counter,progressbar=foot,sectionpage=none]{metropolis}
\usecolortheme[snowy]{owl}
\setbeamercolor{progress bar}{fg=OwlBlue}
\setbeamertemplate{itemize item}{\color{OwlBlue!80}$\blacktriangleright$}
\setbeamertemplate{itemize subitem}{\color{OwlBlue!70}\textbullet}
\setbeamertemplate{enumerate items}[square]
\setbeamercolor{itemize items}{bg=OwlBlue}

\definecolor{dolphin_db}{HTML}{4747ba}
\definecolor{dolphin_lb}{HTML}{aeaee1}
\definecolor{dolphin_b}{HTML}{8585d1}
\setbeamercolor{frametitle}{bg=OwlBlue!75}

\usepackage[utf8]{inputenc}
\usepackage[T1]{fontenc}
\usepackage[english]{babel}
\usepackage{tikz}
\usepackage{appendixnumberbeamer}
\usepackage{amsmath}
\usepackage{amsfonts}
\usepackage{amssymb}
\usepackage{stmaryrd}
\usepackage{mathtools}
\usepackage{bm}
\usepackage{xcolor}
\colorlet{new}{OwlCyan}
\usetikzlibrary{arrows,shapes,positioning,fit}
\usepackage{array}
\usepackage{layouts}
\usepackage[scaled=0.80]{beramono}  % our monospace font
\usepackage{pifont}
\newcommand{\cmark}{\ding{51}}%
\newcommand{\xmark}{\ding{55}}%

\usepackage[backend=biber,style=alphabetic,citestyle=authortitle,abbreviate=true]{biblatex}
%%\addbibresource{rustre.bib}

\DeclareCiteCommand{\footpartcite}[\mkbibfootnote]
  {\usebibmacro{prenote}}
  {\usebibmacro{citeindex}%
   \printnames{labelname}%
   \setunit{\labelnamepunct}%
   \printfield[citetitle]{title}%
   \newunit%
   \printfield{booktitle}%
 }
  {\addsemicolon\space}
  {\usebibmacro{postnote}}

\DeclareMultiCiteCommand{\footpartcites}[\mkbibfootnote]{\footpartcite}{\addsemicolon\space}

\author{Lucas Baudin \and Martin Pépin \and Raphaël Monat \and Auguste Olivry}
\institute[]{}
\title{De Lustre à Rust : Rustre}
\date{}

\newcommand{\inline}[1]{{\ttfamily #1}}
\newcommand{\minline}[1]{\text{\ttfamily #1}}
\newcommand\comment[1]{\hfill \colorbox{blue!70}{#1}}

\usepackage{listings}


\lstdefinelanguage{minils}%
{morekeywords={and,const,else,every,fby,if,in,merge,mod,node,not,or,pre,then,true,type,var,when,with,True,False},%
  morekeywords=[3]{bool,int,real},%
  otherkeywords={=,=>,<=,->},%
  sensitive,%
  morecomment=[l]//,%
  morecomment=[n]{/*}{*/},%
  morestring=[b]",%
  morestring=[b]',%
  morestring=[b]""",%
}[keywords,comments,strings]%

%% From https://github.com/RanaExMachina/beamer-template/blob/483c81e1ac1cabfe4b5f66cc174c5e0dca402944/listing-languages.tex
\lstdefinelanguage{Rust}
{
  alsoletter={=>!&->},
  morekeywords={
          as, break, const, continue, crate, else, enum, extern, false, fn, for,
          if, impl, in, let, loop, match, mod, move, mut, pub, ref, return,
          Self, self, static, struct, super, trait, true, type, unsafe, use,
          where, while,
          abstract, alignof, become, box, do, final, macro, offsetof, override,
          priv, proc, pure, sizeof, typeof, unsized, virtual, yield,
          bool, char, i8, i16, i32, i65, u8, u16, u32, u64, isize, usize, f32,
          f64, str, String
          =>, &self, enum,
  },
  % macros
  emph={println!,vec!,macro_rules!,assert_eq!},
  sensitive=true, % keywords are case-sensitive
  morecomment=[l]{//}, % l is for line comment
  morecomment=[n]{/*}{*/},%
  morecomment=[l]{///}, % l is for line comment
  morecomment=[l]{//!}, % l is for line comment
  morestring=[b]" % defines that strings are enclosed in double quotes
}

\definecolor{mygreen}{rgb}{0,0.6,0}
\definecolor{mygray}{rgb}{0.5,0.5,0.5}
\definecolor{mymauve}{rgb}{0.58,0,0.82}

\lstset{ %
  %backgroundcolor=\color{white},   % choose the background color; you must add \usepackage{color} or \usepackage{xcolor}
  basicstyle=\ttfamily\footnotesize,        % the size of the fonts that are used for the code
  breakatwhitespace=false,         % sets if automatic breaks should only happen at whitespace
  breaklines=true,                 % sets automatic line breaking
%  captionpos=b,                    % sets the caption-position to bottom
  commentstyle=\color{mygreen},    % comment style
  deletekeywords={...},            % if you want to delete keywords from the given language
  escapechar=!,
  escapeinside={\%*}{*)},          % if you want to add LaTeX within your code
  extendedchars=true,              % lets you use non-ASCII characters; for 8-bits encodings only, does not work with UTF-8
  keepspaces=true,                 % keeps spaces in text, useful for keeping indentation of code (possibly needs columns=flexible)
  keywordstyle=\color{blue},       % keyword style
  otherkeywords={*,...},           % if you want to add more keywords to the set
  numbers=none,                    % where to put the line-numbers; possible values are (none, left, right)
  numbersep=8pt,                   % how far the line-numbers are from the code
  numberstyle=\tiny\color{mygray}, % the style that is used for the line-numbers
  %rulecolor=\color{black},         % if not set, the frame-color may be changed on line-breaks within not-black text (e.g. comments (green here))
  showspaces=false,                % show spaces everywhere adding particular underscores; it overrides 'showstringspaces'
  showstringspaces=false,          % underline spaces within strings only
  showtabs=false,                  % show tabs within strings adding particular underscores
%  stepnumber=5,  % the step between two line-numbers. If it's 1, each line will be numbered
  firstnumber=1,numberfirstline=true,
  stringstyle=\color{mymauve},     % string literal style
  tabsize=2,	                   % sets default tabsize to 2 spaces
  %title=\lstname,                   % show the filename of files included with \lstinputlisting; also try caption instead of title
%  frame=Ltb,
  language=minils,
  keywordsprefix=\#,
  alsoletter=\#,
%  breaklines=true
%  literate={\#}{{/}}
}



\newcommand\blfootnote[1]{%
  \begingroup
  \renewcommand\thefootnote{}\footnote{#1}%
  \addtocounter{footnote}{-1}%
  \endgroup
}


\begin{document}

\maketitle

\section{\null}
\subsection{Introduction}

\begin{frame}{\null}
  \begin{tikzpicture}[node distance=1cm, align=center]
    \node (parsing)         [rectangle, rounded corners, draw=black, fill=blue!50, minimum height = .75cm, minimum width=3cm]                            {Parsing};
    \node (typing)          [rectangle, rounded corners, draw=black, fill=blue!50, minimum height = .75cm, minimum width=3cm, right = of parsing]        {Typing};
    \node (clocking)        [rectangle, rounded corners, draw=black, fill=blue!50, minimum height = .75cm, minimum width=3cm, right = of typing]         {Clocking};
    \node (normalization)   [rectangle, rounded corners, draw=black, fill=blue!50, minimum height = .75cm, minimum width=3cm, below = of clocking]       {Normalization};
    \node (scheduling)      [rectangle, rounded corners, draw=black, fill=blue!50, minimum height = .75cm, minimum width=3cm, left = of normalization]   {Scheduling};
    \node (object)          [rectangle, rounded corners, draw=black, fill=blue!50, minimum height = .75cm, minimum width=3cm, left = of scheduling]      {Object generation};
    \node (extraction)      [rectangle, rounded corners, draw=black, fill=blue!50, minimum height = .75cm, minimum width=3cm, below = of object]         {Code extraction};
    \node (rust)            [right = of extraction] {\includegraphics[width=0.1\textwidth]{img/rust.png}};
    \node (why3)            [below = of rust] {Why3};

    \path[draw, ->, line width = .3mm] (parsing) -- (typing);
    \path[draw, ->, line width = .3mm] (typing) -- (clocking);
    \path[draw, ->, line width = .3mm] (clocking) -- (normalization);
    \path[draw, ->, line width = .3mm] (normalization) -- %%node[yshift=.5cm] {\footnotesize Object AST}
                    (scheduling);
    \path[draw, ->, line width = .3mm] (scheduling) -- (object);
    \path[draw, ->, line width = .3mm] (object) -- (extraction);
    \path[draw, ->, line width = .3mm] (extraction) -- (rust);
    \path[draw, ->, line width = .3mm] (extraction.south) |- (why3.west);
  \end{tikzpicture}
\end{frame}

\subsection{Pong}

\begin{frame}{\null}
  \centering \Large Pong : démo !
\end{frame}

\subsection{Typage - GADTs}
\begin{frame}[fragile]{\null}
  \begin{lstlisting}[language=caml]
type enum

type 'a num_ty =
  | TyZ : int num_ty
  | TyReal : float num_ty

type _ ty =
  | TyBool : bool ty
  | TyNum  : 'a num_ty -> 'a num_ty ty
  | TyEnum : string * string list -> enum ty

type 'a expr_desc =
  | EConst : 'a const -> 'a ty expr_desc
  | EFby   : 'a const * 'a ty expr -> 'a ty expr_desc
  | EBOp   : ('a, 'b) binop * 'a ty expr * 'a ty expr -> 'b ty expr_desc
  ...
  \end{lstlisting}
\end{frame}

\subsection{Horloges}
\begin{frame}[fragile]{-- Inférence}
  \begin{itemize}
  \item Algo W avec effets de bords, références...
  \end{itemize}
\end{frame}

\begin{frame}[fragile]{-- Exemple}
  \begin{columns}
    \begin{column}{0.5\textwidth}
      \begin{lstlisting}[language=minils]
node sum(i: int) = (s: int)
with
     s = i -> pre s + i

node main (x:bool) = (t1, t2:int)
with var t1, t2:int; half:bool in
     half = true -> not (pre half);
     t1 = sum(1 when true(half));
     t2 = sum(1) when true(half)
   \end{lstlisting}
 \end{column}
 \begin{column}{0.5\textwidth}
   \pause
   \begin{lstlisting}[escapechar=!]
node sum
    i : [!$\alpha$!]
    s : [!$\alpha$!]

node main
    x : [base]
    half : [!$\alpha$!]
    t1 : [!$\alpha$! on True(half)]
    t2 : [!$\alpha$! on True(half)]
  \end{lstlisting}
 \end{column}
\end{columns}
\pause
\begin{center}
  \begin{tabular}{|c|ccccccc|}
    \hline
    t1    & 1 &  & 2 & & 3 & & 4\\
    \hline
%%    sum(1)& 1 & 2 & 3 & 4 & 5 & 6 & 7\\
%%    \hline
    t2    & 1 &  & 3 & & 5 & & 7\\
    \hline
  \end{tabular}
\end{center}
\end{frame}


\subsection{Extraction Rust}

\begin{frame}[fragile]{\null}
  \begin{itemize}
  \item<1-> Un noeud = un module
  \item<2-> Variables immuables \onslide<3->{(cas des structures de contrôle)}
  \item<4-> \inline{\#[derive(Default)]}
  \end{itemize}
\end{frame}

\begin{frame}[fragile]{\null}
  \begin{tabular}{c}
    \begin{lstlisting}[language=minils]
node sqrt (s:real) = (res:real)
with var x:bool in
  res = 0.5*(s+1) -> (0.5 * (pre res + (s / pre res)))

node main (s:real) = (res:real)
with
  res = sqrt(s) every (false -> (s <> pre s))
\end{lstlisting}
  \end{tabular}
  
  \begin{center}
    \begin{tabular}{|c|ccccccc|}
      \hline
      Input  & 16   & 16   & 16   & 16   & 2   & 2    & 2   \\
      \hline
      Output & 8.5  & 5.19 & 4.14 & 4.00 & 1.5 & 1.42 & 1.41 \\
      \hline
    \end{tabular}
  \end{center}
\end{frame}

\begin{frame}[fragile]{\null}
  %% + main = loop { parse, give to step, print result }
  \begin{lstlisting}[language=Rust]
mod node_sqrt {
  pub struct Machine { ... }

  impl Machine {
    pub fn step(&mut self, ...) -> ... { ... }
    pub fn reset(&mut self) { ... }
  }
}

mod node_main {
  use node_sqrt;

  pub struct Machine { ... }
  impl Machine {
    pub fn step(&mut self, ...) -> ... { ... }
    pub fn reset(&mut self) { ... }
  }
}
  \end{lstlisting}
\end{frame}

\begin{frame}[fragile]{\null}
  \begin{lstlisting}[language=Rust]
mod node_main {
  use node_sqrt;

  #[derive(Default)]
  pub struct Machine { var3:f64, inst1:node_sqrt::Machine }

  impl Machine {
    pub fn step(&mut self, s:f64) -> f64 {
      // ...
      let res = self.inst1.step(s);
      // ...
    }

    pub fn reset(&mut self) {
      self.inst1.reset(); self.var3 = Default::default();
    }
  }
}
\end{lstlisting}
\end{frame}

\begin{frame}[fragile]{\null}
  \begin{lstlisting}[language=Rust]
pub fn main() {
  let mut mach:node_main::Machine = Default::default();
  loop {
    let arg0 = parse_args();
    println!("{:?}", mach.step(arg0));
  }
}
  \end{lstlisting}
\end{frame}



\subsection{Optimisations}
\begin{frame}{\null}
  \begin{itemize}
  \item Fusion des merges
  \item Simplification des merges triviaux
  \end{itemize}
\end{frame}

\subsection{Extraction Why3}
\begin{frame}{\null}
  \begin{itemize}
    \item préservation sémantique
  \end{itemize}
\end{frame}



\subsection{Conclusion}
\begin{frame}{\null}
  \begin{itemize}
  \item todo
  \end{itemize}
\end{frame}

\appendix
\subsection{Besoin d'une annexe ?}
\begin{frame}{\null}
\end{frame}


\end{document}
